\chapter{绪论}

\section{研究背景与意义}

全球能源结构正经历以化石能源向可再生能源过渡为主线的深刻变革。截至2024年，风力发电和光伏发电在全球新增装机容量中的占比已超过80\%，中国的风光并网装机容量连续多年位居世界首位。随着风光渗透率的持续攀升，电力系统的运行特征发生了根本性改变：传统上由负荷需求单一驱动的调度模式，转变为必须同时应对负荷波动与风光出力不确定性的双重挑战\cite{hong2016probabilistic,kaur2016net}。

在高比例可再生能源并网背景下，净负荷（net load）——即总用电需求减去风电出力与光伏出力之和——成为电力系统调度的核心预测对象。净负荷同时继承了人类用电行为的复杂周期性和天气驱动的随机波动性：光伏受辐照度和云量影响呈现日内"鸭子曲线"特征，风电受风速变化呈现跨日尺度的非平稳波动\cite{hanifi2020critical,inman2013solar}。两者的叠加使得净负荷的多步预测直接关系到机组组合优化、备用容量调度和电力市场出清等关键运营决策。准确预测净负荷对可再生能源高效消纳和电网安全稳定运行具有重要的工程意义\cite{kaur2016net}。

近年来，深度学习方法在风光耦合负荷预测精度方面取得了显著进展。LSTM能够捕捉时间序列中的长程依赖关系\cite{silva2024lstm}；Transformer及其变体通过自注意力机制建模全局时间模式\cite{ebrahimzadeh2024ai}；CNN及其与RNN的混合架构进一步提升了多尺度特征提取能力\cite{wang2023hybrid}。这些方法在标准评估指标上已经大幅超越传统统计模型。然而，深度学习模型本质上是"黑箱"：其内部包含数万至数百万个参数的非线性映射，无法提供可审计的推理路径。电力系统运营者不仅需要准确的预测数值，更希望理解风速、辐照度、温度、历史负荷等关键因素如何驱动预测结果。这一需求在电网安全运行和调度决策的可追溯性方面尤为迫切。

可解释人工智能（explainable AI, XAI）领域的现有方法——如SHAP值和注意力权重可视化——能够提供近似的事后解释，但它们本质上是对黑箱模型行为的局部线性近似或统计归因，无法给出确定性的数学公式\cite{machlev2022explainable}。与之相比，符号回归（symbolic regression）旨在从数据中直接发现闭式数学表达式，是模型可解释性的一种更强形式。然而，传统符号回归方法（如PySR\cite{cranmer2023pysr}）在高维输入场景下面临搜索空间组合爆炸的困难。

Kolmogorov-Arnold网络（Kolmogorov-Arnold Networks, KAN）的提出为这一矛盾提供了新的技术路径。KAN于2025年发表于ICLR\cite{liu2025kan}，其核心思想是将传统MLP中节点上的固定激活函数替换为网络边上的可学习B样条函数。这一架构设计使KAN同时具备两个关键优势：首先，B样条函数的高度灵活性使KAN在较小的参数规模下即可达到接近甚至超过同等规模MLP的预测精度；其次，训练完成后可通过稀疏化正则、结构化剪枝和逐边符号拟合等步骤，将学到的数值映射转换为显式的SymPy数学表达式\cite{liu2025kan,liu2025kan2}。KAN的这种"先学习、后提取"范式在物理学方程重发现任务上已取得成功——Liu和Tegmark\cite{liu2025kan2}在Physical Review X上展示了KAN从数值数据中恢复出Kolmogorov湍流定标律等经典物理方程的能力。

然而，将KAN的符号提取能力从受控的物理方程发现迁移到实际的时间序列预测任务时，出现了一个此前未被系统研究的结构性问题：在包含自回归滞后特征的预测场景中，滞后项（如前一时刻的负荷值、前一时刻的风电出力等）通常解释了短期预测方差的绝大部分。当KAN执行L1+熵正则化驱动的稀疏化时，这些高方差贡献的滞后特征形成"捷径"（shortcut），物理气象变量（辐照度、风速、温度等）在稀疏化竞争中被剪除。最终提取的符号公式退化为仅含负荷滞后项的极简表达，全部物理含义丢失。这一现象不是KAN或符号回归方法本身的缺陷，而是高频时间序列中自回归特征与外生物理特征之间信息竞争的结构性后果。理解这一机制，并找到原则性的应对方案，是KAN符号回归在能源预测领域实际应用的前提条件。

Chen等\cite{chen2025mckan}提出的MCKAN（多尺度卷积Kolmogorov-Arnold网络）在风光功率多步预测中取得了优异精度，证明了KAN架构在能源预测任务上的有效性。但MCKAN的研究重心在预测性能优化，并未涉及符号提取和公式可解释性。本课题从MCKAN的多步预测框架出发，聚焦于从KAN到符号公式的提取过程中物理变量的可辨识性问题，构成对现有工作在可解释性维度上的补充。

\section{研究现状与不足}

围绕本课题的核心问题——KAN符号提取中物理变量为何消失——可沿三条技术脉络梳理研究现状。

\textbf{风光耦合负荷预测方面。}负荷预测经历了从统计方法（ARIMA、指数平滑）到机器学习（随机森林、支持向量机）再到深度学习的演进路径\cite{hong2016probabilistic,ebrahimzadeh2024ai}。在净负荷预测这一细分方向上，Kaur等\cite{kaur2016net}较早地指出风光出力的不确定性叠加是净负荷预测的核心难点。Silva-Rodriguez等\cite{silva2024lstm}将LSTM应用于含风光微电网的净负荷预测，验证了深度序列模型在该任务上的有效性。Ebrahimzadeh等\cite{ebrahimzadeh2024ai}在对短期负荷预测AI技术的全面综述中明确指出，精度与可解释性之间的矛盾是该领域的核心瓶颈——现有高精度方法几乎无一能输出可审计的数学公式。

\textbf{KAN与符号提取方面。}KAN\cite{liu2025kan}在ICLR 2025上的发表引发了广泛关注。其关键突破在于证明了边上的可学习激活函数（B样条）在训练后可通过稀疏化-剪枝-符号拟合的链路转换为解析公式。Liu和Tegmark\cite{liu2025kan2}在PRX上系统展示了KAN在物理方程发现中的符号回归能力，成功恢复了包括Stokes定律、相对论时间膨胀等12个经典物理方程。Panczyk等\cite{panczyk2024kan}在核工程领域验证了KAN-to-symbol的转换路径，进一步扩展了KAN符号提取的应用范围。然而，这些成功案例有一个共同特征：输入变量数目有限（通常2--5个），且不包含自回归滞后特征。在这种条件下，所有输入变量都承载独立的物理信息，符号提取的稀疏化过程不存在结构性的信息竞争。

\textbf{KAN在能源预测中的应用方面。}Chen等\cite{chen2025mckan}的MCKAN采用多尺度卷积与KAN的组合架构，在15--45分钟尺度的风光功率预测中达到了优于Transformer等基线的精度。该工作采用Pearson相关系数进行特征选择、Min-Max归一化进行数据预处理，其多步预测设计和子任务分解思路对本课题具有方法论层面的参考价值。但MCKAN聚焦于预测性能的提升，并未执行从KAN到符号公式的提取步骤，也未考察物理变量在稀疏化过程中的命运。Mubarak等\cite{mubarak2024kan}将KAN应用于风电场功率预测，同样以预测精度为主要评价维度，未涉及符号可解释性。

\textbf{符号回归方面。}基于遗传编程的符号回归方法——从Eureqa\cite{schmidt2009distilling}到PySR\cite{cranmer2023pysr}和Operon\cite{burlacu2020operon}——构成了方程发现的主流技术路线。近年来，基于神经网络的符号回归也取得了进展：AI Feynman\cite{udrescu2020ai}利用物理对称性先验引导搜索；基于Transformer的方法\cite{biggio2021neural,kamienny2022end}将方程生成建模为序列翻译任务；SR-LLM\cite{shojaee2025sr}进一步引入大语言模型的检索增强生成能力。SRBench系列基准测试\cite{lacava2021srbench,orzechowski2018where,defranca2023srbench}为符号回归方法的比较提供了标准化平台。但现有符号回归研究——无论是传统GP方法还是神经方法——均面向静态回归问题（给定输入-输出对，发现映射关系），尚未系统考察包含自回归特征的时间序列场景中的符号可辨识性问题。

\textbf{核心研究缺口。}综合以上三条脉络，当前文献存在一个明确的空白：没有研究从机制层面分析KAN符号提取中物理变量被剪除的原因。现有KAN科学应用\cite{liu2025kan,liu2025kan2,panczyk2024kan}工作于干净的物理场景，不存在自回归竞争问题；现有KAN能源预测\cite{chen2025mckan,mubarak2024kan}关注精度但不做符号提取；现有符号回归工作\cite{schmidt2009distilling,cranmer2023pysr,udrescu2020ai,shojaee2025sr}面向静态问题但不处理时间序列的自回归结构。本课题正是针对这一交叉空白展开研究：在KAN-to-symbol的提取管线中，自回归滞后特征的捷径竞争是否是物理变量被稀疏化剪除的主要机制？如何通过可控的干预实验验证这一机制？如何通过结构化分解恢复物理变量的可辨识性？

\section{研究目标与内容}

本课题的研究思路并非一开始就直接指向"捷径竞争"这一机制解释。项目早期首先比较KAN与PySR等方法的预测能力，但需要区分历史探索期与当前主实验期：历史探索期曾出现个别弱配置run中PySR优于KAN的现象；进入当前canonical matched $\Delta\text{net\_load}_{h6}$主实验设置后，KAN teacher在预测层面已经稳定优于direct PySR。由此，问题焦点从"KAN能不能赢"转向"高精度KAN teacher能否转化为保留物理量的符号公式"。

接下来的研究收敛过程是：先在direct路线下观察到公式退化为负荷主导；随后通过focused teacher与component teacher式的定向实验验证，并不是所有物理量都无法进入公式，其中solar构成稳定正例；最终将异常收敛到wind这一边界案例，其更常见的问题不是"训练失败"，而是原始风速项更难以被简洁、稳定、泛化良好的符号表达显式保留，且其可辨识性随horizon呈非单调变化。本文所称teacher均指``KAN teacher + symbolic extraction''的后验提取链路，不对应经典教师-学生蒸馏。

针对上述研究缺口，本课题设定以下四个研究目标：

\textbf{目标一：构建KAN教师-符号提取管线并验证其预测有效性。}在ARPA-E PERFORM公开数据集上，构建包含数据准备、KAN训练（预热-稀疏化-剪枝）、符号提取和多基线对比的完整管线。在净负荷变化量预测任务上，比较KAN教师模型与参数对齐MLP、特征对齐PySR和持续性基线的预测技能分数（skill score），验证KAN教师作为符号提取起点的预测能力。

\textbf{目标二：提出面向符号可辨识性的量化诊断指标体系。}定义变量进入率（Variable Entry Rate, VER）、公式出现率（Formula Appearance Rate, FAR）、活跃边数（active edge count）和迁移间隙比（Transfer Gap Ratio, TGR）四个指标，分别从稀疏结构层面和符号公式层面量化物理变量的可辨识性。该指标体系为回答"哪些物理量能进入公式、哪些被剪除"提供可度量的判据。

\textbf{目标三：设计自回归阻断干预实验，检验捷径竞争机制。}提出S2阻断实验（S2 autoregressive blocking）：在保持所有其他训练设置不变的条件下，仅移除自回归滞后特征，观察物理变量的VER变化。通过$\Delta\VER$（阻断后VER减去阻断前VER）及其bootstrap置信区间，以干预而非观察的方式检验“自回归捷径竞争是否构成物理量被剪除的主导解释”这一假说。

\textbf{目标四：设计S3结构化分解方案作为建设性解决路径。}将净负荷预测分解为负荷、风电和光伏三个子任务，每个子任务使用与其物理本质对应的特征子集训练独立的KAN教师并提取符号公式。通过固定加法组合规则（净负荷 = 负荷 $-$ 风电 $-$ 光伏）合成结构化组合预测。验证该分解策略在保持预测精度的同时，是否能使各子公式稳定保留目标物理变量。

围绕上述四个目标，本课题的具体研究内容包括：

（1）数据管线的设计与实现。采用ARPA-E PERFORM数据集的5分钟采样电力数据和小时级气象观测数据，构造包含历史滞后特征、气象/太阳位置协变量和周期编码的26维输入特征。采用时间顺序划分策略，并在各子集间设置防泄漏间隔。

（2）KAN训练与符号提取管线。实现预热-稀疏化（L1+熵正则化）-剪枝-符号拟合的四阶段训练流程。对剪枝后存活的网络边逐一执行候选函数拟合，经SymPy化简生成闭式公式。

（3）公式坍缩现象的系统记录。在直接净负荷目标上，以多组符号配置（3种候选函数库$\times$3种拟合阈值）执行符号提取，记录所有配置下公式坍缩为仅含负荷项的现象。

（4）S2阻断干预实验。在聚焦风电任务和直接净负荷任务上，分别执行5个随机种子的阻断/非阻断对比实验，计算$\Delta\VER$及其95\% bootstrap置信区间。

（5）S3结构化分解与组合评估。训练负荷、风电、光伏三个子任务的KAN教师模型，提取各自的符号公式，合成结构化组合预测并与直接KAN教师对比精度和可解释性。

\section{论文组织结构}

本文共分七章，各章内容安排如下：

第一章为绪论，阐述了风光耦合负荷预测中可解释性需求的工程背景，指出KAN符号提取面临的自回归捷径竞争问题，明确了研究目标和具体研究内容。

第二章为文献综述，系统梳理了负荷与净负荷预测方法、符号回归与方程发现、Kolmogorov-Arnold网络、MCKAN与风光预测、物理信息学习与可解释性等方向的研究进展，归纳了现有工作的不足和本课题的切入点。

第三章为数据与问题定义，介绍ARPA-E PERFORM数据集的结构与预处理方法，定义净负荷和差分目标的计算方式，说明特征工程和数据划分策略。

第四章为方法设计，详细描述KAN教师训练管线、符号提取流程、VER/FAR/TGR等可辨识性诊断指标的定义，S2阻断干预实验的对照协议，以及S3结构化分解方案的设计。

第五章为实验结果与分析，按四幕结构展开：预测精度基线（Act 1）、公式坍缩现象（Act 2）、S2阻断干预结果（Act 3）和S3分解评估（Act 4），以$\Delta\VER$和bootstrap置信区间为核心证据链。

第六章为讨论，从捷径竞争机制的普遍性、实验过程的收敛链条、精度-可解释性权衡、方法局限性（数据泄漏风险、种子数量、单一数据集）等角度展开分析，并明确论文叙事边界。

第七章为结论与后续工作，总结本课题的主要贡献和发现，并说明总公式收口工作与后续研究方向。
